% ============================================================================
%  El kernel Gabor como atomo: envolvente x onda
%  Documento tecnico optimizado para lectura en Kindle Scribe (10.2", e-ink).
%
%  Compilar:  pdflatex docs/kernel_gabor_atomo.tex   (dos veces, por el indice)
%
%  DECISIONES DE FORMATO PARA E-INK (no tocar sin motivo):
%   - papersize 6.1x8.15 in = la relacion 3:4 del Scribe -> el PDF llena la
%     pantalla sin zoom ni scroll horizontal, que es lo que arruina la lectura.
%   - margenes 0.32 in: en un lector no hay que sujetar el papel con el pulgar.
%   - cuerpo 12pt sobre pagina pequena = tipo grande en pantalla.
%   - XCharter: serif de trazo grueso, pensada para tinta; en e-ink las
%     modernas de trazo fino (Computer Modern) se ven lavadas.
%   - sin color: el Scribe es escala de grises. Los enlaces van en negro.
% ============================================================================
\documentclass[12pt]{extarticle}

\usepackage[T1]{fontenc}
\usepackage[utf8]{inputenc}
\usepackage[spanish,es-noquoting]{babel}
\usepackage{XCharter}
\usepackage[scaled=0.96]{helvet}
\usepackage{inconsolata}
\usepackage{microtype}
\usepackage{amsmath,amssymb}
\usepackage{booktabs}
\usepackage{array}
\usepackage{listings}
\usepackage{enumitem}
\usepackage{fancyhdr}
\usepackage[
  paperwidth=6.1in, paperheight=8.15in,
  top=0.40in, bottom=0.46in, left=0.32in, right=0.32in,
  footskip=0.24in
]{geometry}
\usepackage{float}
\usepackage{pgfplots}
\pgfplotsset{compat=1.16}
\usepackage[hidelinks]{hyperref}

% --- Estilo unico para TODAS las figuras (e-ink: sin color, se distingue por
%     trazo y grosor, nunca por tono) -----------------------------------------
\pgfplotsset{
  eink/.style={
    width=12.4cm, height=5.3cm,
    scale only axis,
    axis lines=left,
    tick align=outside,
    tick style={black, thin},
    every axis label/.append style={font=\small},
    tick label style={font=\footnotesize},
    legend style={font=\footnotesize, draw=none, fill=none,
                  at={(0.5,-0.26)}, anchor=north, legend columns=2,
                  cells={anchor=west}, inner sep=1.5pt,
                  /tikz/every even column/.append style={column sep=0.45cm}},
    label style={font=\small},
    line width=0.9pt,
    samples=180,
  },
  /pgf/number format/use comma,
}
\newcommand{\figcap}[1]{\caption{\small #1}}

% --- Espaciado pensado para pantalla: sin sangria, aire entre parrafos -------
\setlength{\parindent}{0pt}
\setlength{\parskip}{0.55em}
\linespread{1.06}
\setlist{leftmargin=1.15em, itemsep=0.20em, topsep=0.30em, parsep=0pt}

\sloppy
\emergencystretch=2.2em

\pagestyle{fancy}
\fancyhf{}
\fancyfoot[C]{\footnotesize\thepage}
\renewcommand{\headrulewidth}{0pt}

\lstset{
  basicstyle=\ttfamily\footnotesize,
  columns=fullflexible,
  keepspaces=true,
  breaklines=true,
  breakindent=0pt,
  frame=leftline,
  framesep=5pt,
  xleftmargin=7pt,
  rulecolor=,
  aboveskip=0.7em, belowskip=0.7em,
  showstringspaces=false,
}

% Notacion
\newcommand{\ssum}{\Sigma a}
\newcommand{\code}[1]{\texttt{#1}}

\title{\bfseries El kernel Gabor como átomo\\[0.25em]
\large envolvente $\times$ onda: qué es, y qué le hace\\ a cada fase de la optimización}
\author{\normalsize Proyecto 2DGS $\rightarrow$ Beta / Gabor Splatting}
\date{\normalsize 31 de julio de 2026}

\begin{document}
\maketitle
\thispagestyle{fancy}

{\small
Este documento describe el kernel implementado en la rama
\code{gabor-atomo-envolvente-onda} (commit \code{54274f5} y posteriores) y sigue
su rastro por todas las fases del entrenamiento: preprocesado, forward,
backward, proyección en Python, densificación y serialización. Cada fórmula
corresponde a una línea concreta del código, citada como
\code{archivo:línea}. Las cifras experimentales son las de los runs
5, 6 y 7 sobre \emph{flowers}.
}

\tableofcontents
\clearpage

% ===========================================================================
\section{El problema que resuelve}

\subsection{Dos palabras que se repiten en todo el documento}

\textbf{Tent.} Del inglés \emph{tienda de campaña}. Es el nombre corto del
perfil
\[
E(r)=(1-r)^{\beta},
\]
que vale $1$ en el centro del surfel y baja hasta $0$ en el borde. El nombre
viene del caso $\beta=1$, donde el perfil es literalmente un triángulo --- en
2D, un cono, o sea una tienda de campaña. Con otros $\beta$ ya no es recto
(se comba hacia dentro o hacia fuera, figura~\ref{fig:envolvente}), pero el
nombre se ha quedado.

Importa por dos motivos que van más allá de la geometría:

\begin{itemize}
\item Es el \textbf{kernel de referencia} del proyecto. El run\,67 --- el mejor
      resultado del camino clásico, $20{,}6684$ de PSNR --- usa exactamente
      este perfil, fijo, sin nada que aprender. Es contra él contra quien hay
      que ganar, no contra el 2DGS oficial.
\item Es el \textbf{punto neutro} del kernel nuevo: con los coeficientes a
      cero, el átomo \emph{es} la tent. Por eso, cuando en el documento se lee
      ``ese splat es un tent'', significa que su parte aprendida no está
      haciendo nada.
\end{itemize}

\textbf{Percentil.} El percentil $p_{XX}$ de un conjunto de valores es el valor
que deja por debajo al $XX$\,\% de ellos. El $p_{50}$ es la mediana: la mitad
de los splats está por debajo y la mitad por encima. El $p_{90}$ deja debajo al
90\,\%, así que solo uno de cada diez lo supera.

En este proyecto los diagnósticos se leen \textbf{siempre} en percentiles
($p_{10}$, $p_{50}$, $p_{90}$) y nunca en media o máximo, y es una regla
aprendida a base de errores:

\begin{itemize}
\item El \textbf{máximo} de una nube de $4{,}9$ millones de splats es un solo
      splat. Puede ser enorme y no significar nada. Varias veces se dio por
      activo un mecanismo mirando su valor máximo cuando en realidad estaba
      apagado para el 99\,\% del modelo.
\item La \textbf{media} la arrastra la cola. En distribuciones tan asimétricas
      como estas puede quedar por encima del 90\,\% de los datos.
\item El \textbf{histograma} responde a la pregunta que de verdad importa:
      ¿a cuánta nube le está pasando esto? Cuando se lee ``$p_{50}$ de
      $f\!\cdot\!W$ vale $0{,}328$'' hay que entender: \emph{la mitad de los
      splats tiene la onda por debajo de un tercio de lo necesario para
      oscilar}. Eso es un diagnóstico; el máximo no lo habría sido.
\end{itemize}


\subsection{El kernel anterior y su punto de degeneración}

Hasta el run\,4 el kernel de opacidad era una serie de cosenos elevada a
$\beta$:
\[
K_{\text{legacy}}(r)=\Bigl(\frac{f(r)}{f(0)}\Bigr)^{\beta},\qquad
f(r)=\tfrac12+\sum_{n=1}^{3} a_n\cos\bigl((2n-1)\pi r\bigr),
\]
con $r=\sqrt{\rho}$ la distancia normalizada al centro del surfel y
$a_n\ge 0$, $\sum a_n\le\frac12$ impuesto por proyección.

Tiene un defecto estructural que no se ve mirando la fórmula, sino
preguntándose \emph{qué kernel sale cuando el modelo no aprende nada}. Con
$a=0$ queda $f\equiv\frac12$, es decir
\[
K_{\text{legacy}}\big|_{a=0}\equiv 1 \quad\text{en todo el soporte},
\]
un \textbf{disco plano de opacidad uniforme}: el peor kernel posible. Y ese
punto no está en un rincón cualquiera del espacio de parámetros: está
\emph{en la frontera} $a_n\ge 0$, justo donde el optimizador empuja cuando
quiere apagar la modulación.

Peor aún: el kernel de referencia, la \emph{tent} $(1-r)^\beta$ del run\,67,
se obtiene con $a=(4/\pi^2)(1,\tfrac19,\tfrac1{25})$ renormalizado, lo que
consume $\ssum=0{,}4665$, el \textbf{93{,}3\,\%} del presupuesto disponible
$\ssum\le\frac12$. Dicho de otro modo: en el kernel legacy, \emph{limitarse a
reproducir la envolvente ya agota casi todo el presupuesto}, y lo que queda
para modular la forma es una miseria.

\subsection{La observación de fondo}

Un átomo de Gabor es, por definición, un producto de dos factores
independientes: una \textbf{envolvente} que decide hasta dónde llega el
soporte, y una \textbf{onda} que decide cómo oscila dentro de él. El kernel
legacy no tiene esa separación: el decaimiento y la oscilación salen los dos
del mismo $\sum a_n\cos(\cdot)$, así que los tres coeficientes hacen dos
trabajos con un solo presupuesto.

La medida que lo confirma: solo el \textbf{7{,}9\,\%} del conjunto factible
$\{a_n\ge0,\ \ssum\le\frac12\}$ da un perfil decreciente. El otro 92\,\% son
kernels con lóbulos secundarios --- los anillos concéntricos que se veían en
el césped.

\subsection{La solución: factorizar}

\begin{equation}
\boxed{\;K(r,t) \;=\; \underbrace{(1-r)^{\beta}}_{\text{envolvente }E}
\;\cdot\; \underbrace{S(\varphi t)}_{\text{onda}}\;}
\label{eq:kernel}
\end{equation}
\begin{equation}
S(x) \;=\; b \;+\; \sum_{n=1}^{3} a_n\cos\bigl((2n-1)\,x\bigr),
\label{eq:onda}
\end{equation}
\begin{equation}
b \;=\; \gamma + (1-\gamma)\,\bigl(1-\ssum\bigr),\qquad \ssum=\sum_n a_n .
\label{eq:pedestal}
\end{equation}

El término $b$ es el \textbf{pedestal}, tomado de AdaGaR. Es la pieza que lo
cambia todo, porque hace que
\[
a=0 \;\Longrightarrow\; b=1 \;\Longrightarrow\; S\equiv 1
\;\Longrightarrow\; K=(1-r)^{\beta},
\]
\emph{la tent del run\,67 exactamente}. Y el modelo se inicializa ahí
(\code{gaussian\_model.py:298}, \code{a\_init\_row}). Las consecuencias no son
cosméticas:

\begin{itemize}
\item El \textbf{baseline es el punto neutro} del espacio de parámetros, no una
      esquina. Aprender la forma solo puede sumar; si no encuentra nada, se
      queda donde estaba.
\item La envolvente es \textbf{gratis}: ya no se paga con el presupuesto de
      $a$. Los tres coeficientes quedan libres para lo único que deben hacer.
\item La fracción de kernels decrecientes dentro del conjunto factible sube de
      7,9\,\% a \textbf{70,8\,\%} (Monte Carlo convergido).
\item El control experimental es exacto: las primeras iteraciones deben
      reproducir el run\,67 dígito a dígito.
\end{itemize}

% ===========================================================================
\section{Definición formal}

\subsection{Los tres modos}

El modo lo fija la variable de entorno \code{GABOR\_KERNEL} y viaja a CUDA
como el entero \code{gabor\_mode} (\code{gabor\_kernel.h:47--49}):

\begin{center}\small
\begin{tabular}{@{}llp{2.5cm}@{}}
\toprule
\textbf{modo} & \textbf{id} & \textbf{kernel} \\
\midrule
\code{legacy} & 0 & $(f/f_0)^\beta$ \\
\code{radial} & 1 & $(1-r)^\beta\,S(\varphi r)$ \\
\code{dir}    & 2 & $(1-r)^\beta\,S(\varphi u)$ \\
\bottomrule
\end{tabular}
\end{center}

La diferencia entre \code{radial} y \code{dir} es el argumento de la onda:

\begin{itemize}
\item \textbf{radial}: $t=r$, la modulación es concéntrica. Produce
      \emph{anillos}.
\item \textbf{dir}: $t=u$, la primera coordenada local del surfel. La
      modulación son \emph{franjas paralelas} en una dirección, que es lo que
      un átomo de Gabor es de verdad. La orientación de las franjas la hereda
      de la rotación del surfel, así que \textbf{ya es aprendible} sin añadir
      ningún parámetro.
\end{itemize}

Los runs 5, 6 y 7 usan \code{dir}, que en el A/B local ganó con \emph{la mitad
de splats} que \code{radial} (287\,k contra 597\,k): más expresividad por
primitiva, que es la promesa del átomo.

\subsection{La frecuencia: el modo \code{world} y por qué importa}

$\varphi$ es adimensional dentro de CUDA, pero se construye en Python
(\code{gaussian\_model.py:360--371}) de dos formas:

\begin{align}
\text{\code{norm}:}\quad & \varphi=\pi \\[0.2em]
\text{\code{world}:}\quad & \varphi=\frac{f_1}{\text{extent}}\; s_u
\label{eq:phi}
\end{align}

donde $s_u$ es la escala del eje $u$ en unidades de mundo y \emph{extent} es
\code{spatial\_lr\_scale}.

La consecuencia de \eqref{eq:phi} merece detenerse. La coordenada $u$ que
recibe CUDA está \emph{normalizada por la escala del splat} (por eso el
soporte es $\rho<1$), o sea $u = x_{\text{mundo}}/s_u$. Al multiplicar:
\[
\varphi\,u \;=\; \frac{f_1 s_u}{\text{extent}}\cdot\frac{x_{\text{mundo}}}{s_u}
\;=\; \frac{f_1}{\text{extent}}\;x_{\text{mundo}} .
\]
Las escalas \textbf{se cancelan}. La longitud de onda en el mundo es
\[
\lambda = \frac{2\pi\,\text{extent}}{f_1},
\]
\emph{la misma para todos los splats}. El modo \code{world} impone una rejilla
de frecuencia fija en la escena, y cada surfel muestra tantos ciclos como le
quepan según su tamaño. Un splat grande desarrolla textura; uno pequeño no
llega ni a medio ciclo y degenera al tent.

En modo \code{norm} pasa lo contrario: $\varphi=\pi$ significa medio periodo
dentro del soporte \emph{sea cual sea el tamaño}, es decir la frecuencia va
atada al splat y no a la escena.

La figura \ref{fig:worldnorm} enseña la diferencia mejor que las fórmulas.
Toma dos surfels de tamaños muy distintos ($s=0{,}35$ y $s=0{,}10$) y dibuja el
kernel que produce cada modo \emph{en coordenadas de mundo}: no normalizado al
splat, sino como lo vería un píxel de la imagen.

\begin{figure}[H]\centering
\begin{tikzpicture}
\begin{axis}[eink, height=3.9cm,
  title={(a) modo \code{norm}: $\varphi=\pi$},
  title style={font=\small, yshift=-2pt},
  xlabel={}, ylabel={$K$},
  xmin=-0.85, xmax=0.5, ymin=0, ymax=1.2,
  xtick={-0.75,-0.5,-0.25,0,0.25,0.5},
  xticklabels={},
  legend to name=leyendaWN]
\addplot[black!20, thin, forget plot] coordinates {(-0.6283,0) (-0.6283,1.2)};
\addplot[black!20, thin, forget plot] coordinates {(-0.4189,0) (-0.4189,1.2)};
\addplot[black!20, thin, forget plot] coordinates {(-0.2094,0) (-0.2094,1.2)};
\addplot[black!20, thin, forget plot] coordinates {(0,0) (0,1.2)};
\addplot[black!20, thin, forget plot] coordinates {(0.2094,0) (0.2094,1.2)};
\addplot[black!20, thin, forget plot] coordinates {(0.4189,0) (0.4189,1.2)};
\addplot[black, line width=1.1pt, domain=-0.75:-0.05, samples=400]
 {pow(1-abs((x+0.40)/0.35),2.236)*(0.741
  +0.30*cos(deg(8.9760*(x+0.40)))
  +0.05*cos(deg(26.928*(x+0.40)))
  +0.02*cos(deg(44.880*(x+0.40))))};
\addlegendentry{surfel grande ($s=0{,}35$)}
\addplot[black, dashed, domain=0.15:0.35, samples=400]
 {pow(1-abs((x-0.25)/0.10),2.236)*(0.741
  +0.30*cos(deg(31.416*(x-0.25)))
  +0.05*cos(deg(94.248*(x-0.25)))
  +0.02*cos(deg(157.08*(x-0.25))))};
\addlegendentry{surfel pequeño ($s=0{,}10$)}
\end{axis}
\end{tikzpicture}

\begin{tikzpicture}
\begin{axis}[eink, height=3.9cm,
  title={(b) modo \code{world}: $\varphi=f_1 s_u/\text{extent}$},
  title style={font=\small, yshift=-2pt},
  xlabel={posición en el mundo}, ylabel={$K$},
  xmin=-0.85, xmax=0.5, ymin=0, ymax=1.2,
  xtick={-0.75,-0.5,-0.25,0,0.25,0.5},
  xticklabels={{-0,75},{-0,5},{-0,25},0,{0,25},{0,5}},
  legend style={at={(0.5,-0.42)}, anchor=north, legend columns=2}]
\addplot[black!20, thin, forget plot] coordinates {(-0.6283,0) (-0.6283,1.2)};
\addplot[black!20, thin, forget plot] coordinates {(-0.4189,0) (-0.4189,1.2)};
\addplot[black!20, thin, forget plot] coordinates {(-0.2094,0) (-0.2094,1.2)};
\addplot[black!20, thin, forget plot] coordinates {(0,0) (0,1.2)};
\addplot[black!20, thin, forget plot] coordinates {(0.2094,0) (0.2094,1.2)};
\addplot[black!20, thin, forget plot] coordinates {(0.4189,0) (0.4189,1.2)};
\addplot[black, line width=1.1pt, domain=-0.75:-0.05, samples=400]
 {pow(1-abs((x+0.40)/0.35),2.236)*(0.741
  +0.30*cos(deg(30*(x+0.40)))
  +0.05*cos(deg(90*(x+0.40)))
  +0.02*cos(deg(150*(x+0.40))))};
\addlegendentry{surfel grande ($s=0{,}35$)}
\addplot[black, dashed, domain=0.15:0.35, samples=400]
 {pow(1-abs((x-0.25)/0.10),2.236)*(0.741
  +0.30*cos(deg(30*(x-0.25)))
  +0.05*cos(deg(90*(x-0.25)))
  +0.02*cos(deg(150*(x-0.25))))};
\addlegendentry{surfel pequeño ($s=0{,}10$)}
\end{axis}
\end{tikzpicture}
\figcap{Los mismos dos surfels y los mismos coeficientes, dibujados en
coordenadas de mundo. Las líneas grises son una \emph{regla} de paso
$\lambda=2\pi\,\text{extent}/f_1$, idéntica en los dos paneles.
\textbf{(a)~\code{norm}} produce un kernel \emph{autosemejante}: como
$\varphi=\pi$ es medio ciclo del centro al borde \emph{sea cual sea el tamaño},
la onda no llega nunca a oscilar --- queda en una campana, y lo único que
cambia entre los dos surfels es lo ancha que es. Un kernel así no puede
representar una frecuencia \emph{de la escena}, porque su frecuencia la fija el
propio splat: los dos pintan detalle de distinto tamaño aunque estén sobre la
misma pared.
\textbf{(b)~\code{world}} fija $\lambda$ y deja que varíe lo que tiene sentido
que varíe: cuántos ciclos caben. En el grande caben $3{,}3$ y aparecen franjas;
en el pequeño cabe $0{,}95$ y vuelve a ser una campana, es decir, poco más que
un tent. Lo compartido es la longitud de onda, no la fase: cada surfel ancla la
suya en su propio centro.}
\label{fig:worldnorm}
\end{figure}

Y aquí está el germen de la patología que se describe más adelante: en
\code{world}, \textbf{encoger un surfel lo acerca al tent}. La densificación
encoge surfels a miles, así que el modo que hace físicamente correcta la
frecuencia es también el que la apaga cuando el modelo crece.

\subsection{La cota del presupuesto}

El requisito es $S\ge 0$: la onda no debe volverse negativa dentro del
soporte, o el kernel pintaría opacidad negativa (que el forward recorta, pero
recortar es perder gradiente). El peor caso de \eqref{eq:onda} es que los tres
cosenos valgan $-1$ a la vez, así que
\[
S_{\min} = b - \ssum \;=\; \gamma+(1-\gamma)(1-\ssum)-\ssum
\;=\; 1-(2-\gamma)\,\ssum ,
\]
y por tanto
\begin{equation}
\boxed{\;\ssum \;\le\; \frac{1}{2-\gamma}\;}
\label{eq:cota}
\end{equation}
Con $\gamma=0{,}3$ salen $0{,}5882$ (\code{gaussian\_model.py:309--320}). Nótese
que la cota es \emph{más generosa} que el $\frac12$ del legacy: es el mismo
razonamiento, pero teniendo en cuenta que el pedestal sube automáticamente
cuando $\ssum$ baja.

El dial $\kappa$ (\code{GABOR\_KAPPA}) multiplica la cota. Con $\kappa<1$ se
acota además el contraste valle/pico a $(1-\kappa)/(1+\kappa)$. Se dejó en
$1{,}0$ en todos los runs.

\subsection{Cómo se lee el papel de $\gamma$}

$\gamma$ es el suelo del pedestal. Fija cuánta modulación puede llegar a
haber como máximo:
\[
b\in[\gamma,\,1],\qquad b=1 \text{ si } \ssum=0,\qquad
b=\gamma \text{ si } \ssum=1 .
\]
Con $\gamma=0$ el pedestal podría anularse del todo y la envolvente volvería a
depender de $a$; con $\gamma=1$ el pedestal sería constante y la cota
\eqref{eq:cota} se quedaría en $1$. El $0{,}3$ es el valor de AdaGaR.

% ===========================================================================
\section{Atlas del kernel}

Todas las curvas de esta sección son las funciones \emph{reales} del código,
evaluadas con los parámetros medidos en los runs: $\gamma=0{,}3$,
$\beta=2{,}236$ (media del run\,7 a 30\,k) y, cuando se indica, los
coeficientes medios $a=(0{,}1153,\;0{,}1039,\;0{,}1237)$ del run\,7, que dan
$\ssum=0{,}3429$ y $b=0{,}760$.

\subsection{El punto neutro: caja contra tent}

\begin{figure}[H]\centering
\begin{tikzpicture}
\begin{axis}[eink, xlabel={$r$}, ylabel={$K(r)$},
  xmin=0, xmax=1, ymin=0, ymax=1.15]
\addplot[black, dashed, domain=0:1] {1};
\addlegendentry{legacy con $a=0$: caja}
\addplot[black!45, line width=2.4pt, domain=0:1] {pow(1-x,2.236)};
\addlegendentry{tent del run\,67}
\addplot[black, dotted, line width=1.1pt, domain=0:1] {pow(1-x,2.236)};
\addlegendentry{átomo con $a=0$}
\end{axis}
\end{tikzpicture}
\figcap{La figura que resume el rediseño. Cuando el modelo no aprende nada, el
kernel legacy degenera en un \emph{disco plano} y el átomo cae exactamente
sobre la tent: la curva punteada se apoya sobre la gruesa porque son la misma
función. El baseline pasa de ser una esquina del espacio de parámetros a ser su
punto neutro.}
\end{figure}

\subsection{La envolvente y el papel de $\beta$}

\begin{figure}[H]\centering
\begin{tikzpicture}
\begin{axis}[eink, xlabel={$r$}, ylabel={$E(r)=(1-r)^{\beta}$}, legend columns=3,
  xmin=0, xmax=1, ymin=0, ymax=1.02]
\addplot[black, dotted,   domain=0:1] {pow(1-x,0.5)};   \addlegendentry{$\beta=0{,}5$}
\addplot[black, dashed,   domain=0:1] {pow(1-x,1)};     \addlegendentry{$\beta=1$}
\addplot[black, line width=1.6pt, domain=0:1] {pow(1-x,2.236)}; \addlegendentry{$\beta=2{,}24$ (media medida)}
\addplot[black!60, dashdotted, domain=0:1] {pow(1-x,5)};  \addlegendentry{$\beta=5$}
\addplot[black!60, solid, domain=0:1] {pow(1-x,15)};      \addlegendentry{$\beta=15$}
\end{axis}
\end{tikzpicture}
\figcap{$\beta$ es el único parámetro de la envolvente y controla cuánta masa
queda cerca del centro. El soporte es compacto para todo $\beta$: se anula en
$r=1$, que es lo que permite el guard \code{r2>=1 continue} del rasterizer y
lo que hace finito el $\ln(1-r)$ del gradiente. El rango medido en los runs va
de $0{,}096$ a $14$, con media $2{,}24$.}
\label{fig:envolvente}
\end{figure}

\subsection{La onda: dos formas de gastar el mismo presupuesto}

\begin{figure}[H]\centering
\begin{tikzpicture}
\begin{axis}[eink, xlabel={$x=\varphi t$ (rad)}, ylabel={$S(x)$},
  xmin=0, xmax=12.566, ymin=0.3, ymax=1.2,
  xtick={0,3.1416,6.2832,9.4248,12.566},
  xticklabels={$0$,$\pi$,$2\pi$,$3\pi$,$4\pi$}]
\addplot[black, line width=1.4pt, domain=0:12.566]
  {0.76 + 0.2400*cos(deg(x)) + 0.0700*cos(deg(3*x)) + 0.0329*cos(deg(5*x))};
\addlegendentry{espectro decreciente}
\addplot[black, dashed, domain=0:12.566]
  {0.76 + 0.1153*cos(deg(x)) + 0.1039*cos(deg(3*x)) + 0.1237*cos(deg(5*x))};
\addlegendentry{run\,7 medido ($a_3>a_1$)}
\addplot[black!50, dotted, line width=1pt, domain=0:12.566] {0.76};
\addlegendentry{pedestal $b$}
\end{axis}
\end{tikzpicture}
\figcap{Las dos curvas tienen \emph{exactamente} el mismo presupuesto
($\ssum=0{,}343$) y el mismo pedestal, y son kernels completamente distintos.
Arriba, el reparto decreciente da una oscilación limpia de un solo armónico
dominante. Abajo, el reparto que el optimizador eligió de verdad en el run\,7
--- con $a_3>a_1$ --- da una onda dentada, con la energía en el armónico más
alto. Ambas oscilan alrededor de $b$ y ninguna baja de
$S_{\min}=1-(2-\gamma)\ssum=0{,}417$.}
\end{figure}

\subsection{El interruptor: qué significa ``la onda está apagada''}

\begin{figure}[H]\centering
\begin{tikzpicture}
\begin{axis}[eink, xlabel={$u$ (coordenada local, soporte $|u|\le1$)},
  ylabel={$S(\varphi u)$},
  xmin=-1, xmax=1, ymin=0.3, ymax=1.2]
\addplot[black, line width=1.4pt, domain=-1:1]
  {0.76 + 0.1153*cos(deg(0.154*x)) + 0.1039*cos(deg(3*0.154*x)) + 0.1237*cos(deg(5*0.154*x))};
\addlegendentry{$\varphi=0{,}15$ (p10: apagado)}
\addplot[black, dashed, domain=-1:1]
  {0.76 + 0.1153*cos(deg(0.884*x)) + 0.1039*cos(deg(3*0.884*x)) + 0.1237*cos(deg(5*0.884*x))};
\addlegendentry{$\varphi=0{,}88$ (p50: apagado)}
\addplot[black, dotted, line width=1.1pt, domain=-1:1]
  {0.76 + 0.1153*cos(deg(4.214*x)) + 0.1039*cos(deg(3*4.214*x)) + 0.1237*cos(deg(5*4.214*x))};
\addlegendentry{$\varphi=4{,}21$ (p90: encendido)}
\end{axis}
\end{tikzpicture}
\figcap{Las tres curvas usan los \emph{mismos} coeficientes del run\,7 y solo
cambia $\varphi$, tomado de los percentiles 10, 50 y 90 medidos
($f\!\cdot\!W = 0{,}057$, $0{,}328$ y $1{,}564$, con
$\varphi = f\!\cdot\!W/2\sigma_{\text{env}}$). En la mediana la onda no
completa ni un tercio de ciclo dentro del soporte: es una constante ligeramente
inclinada, o sea \emph{un tent}. Esto es lo que cuenta el
diagnóstico \code{[GABOR-W]} cuando dice ``64\,\% apagados'': no que los
coeficientes sean cero, sino que la frecuencia no da para oscilar.}
\end{figure}

\subsection{El kernel completo}

\begin{figure}[H]\centering
\begin{tikzpicture}
\begin{axis}[eink, xlabel={$u$ (corte por $v=0$)}, ylabel={$K=E\cdot S$},
  xmin=-1, xmax=1, ymin=0, ymax=1.2]
\addplot[black!45, line width=2.2pt, domain=-1:1] {pow(1-abs(x),2.236)};
\addlegendentry{tent ($a=0$)}
\addplot[black, dashed, domain=-1:1]
  {pow(1-abs(x),2.236)*(0.76 + 0.1153*cos(deg(0.884*x)) + 0.1039*cos(deg(3*0.884*x)) + 0.1237*cos(deg(5*0.884*x)))};
\addlegendentry{run\,7, $\varphi=0{,}88$ (p50)}
\addplot[black, solid, domain=-1:1]
  {pow(1-abs(x),2.236)*(0.76 + 0.1153*cos(deg(4.214*x)) + 0.1039*cos(deg(3*4.214*x)) + 0.1237*cos(deg(5*4.214*x)))};
\addlegendentry{run\,7, $\varphi=4{,}21$ (p90)}
\end{axis}
\end{tikzpicture}
\figcap{El producto de las dos figuras anteriores, en modo \code{dir} y a lo
largo del eje de modulación. La envolvente impone el soporte compacto y la
onda esculpe franjas dentro de él. Obsérvese el pico central: vale
$1+\gamma\ssum=1{,}103$, \emph{por encima} de la tent, porque los tres cosenos
entran en fase en $u=0$. Y obsérvese la curva de la mediana: sigue siendo un
tent, solo que un 10\,\% más opaco en el centro y algo más apuntado --- la
campana de la figura anterior no añade estructura, reajusta la forma de la
envolvente. Para el 64\,\% de la nube, eso es todo lo que el kernel aprendido
llega a ser.}
\end{figure}

\subsection{El presupuesto y de dónde sale la cota}

\begin{figure}[H]\centering
\begin{tikzpicture}
\begin{axis}[eink, xlabel={$\ssum$}, ylabel={},
  xmin=0, xmax=0.70, ymin=-0.2, ymax=1.25,
  xtick={0,0.1,0.2,0.3,0.4,0.5,0.5882,0.7},
  xticklabels={0,{0,1},{0,2},{0,3},{0,4},{0,5},{0,59},{0,7}},
  legend columns=2]
\addplot[black, line width=1.4pt, domain=0:0.70] {0.3+0.7*(1-x)};
\addlegendentry{pedestal $b$}
\addplot[black, dashed, domain=0:0.70] {1-1.7*x};
\addlegendentry{$S_{\min}=1-(2-\gamma)\ssum$}
\addplot[black, dotted, line width=1.1pt, domain=0:0.70] {1+0.3*x};
\addlegendentry{$S_{\max}=1+\gamma\ssum$}
\addplot[black!55, line width=1pt, dashdotted] coordinates {(0.5882,-0.2) (0.5882,1.25)};
\addlegendentry{cota $1/(2-\gamma)$}
\addplot[black!55, thin] coordinates {(0,0) (0.70,0)};
\end{axis}
\end{tikzpicture}
\figcap{La cota \eqref{eq:cota} no es un número elegido a mano: es
\emph{exactamente} el punto donde $S_{\min}$ toca cero, o sea donde la onda
empezaría a pintar opacidad negativa. A la izquierda de la línea vertical todo
es válido. Nótese que $S_{\max}$ crece con $\ssum$ mientras $b$ decrece: el
pedestal cede terreno para dejar sitio a la modulación, y por eso la envolvente
sale gratis. Los valores medios medidos ($\ssum\approx0{,}34$) están a poco más
de la mitad del presupuesto.}
\end{figure}

\subsection{La anchura efectiva, o por qué el radio engaña}

\begin{figure}[H]\centering
\begin{tikzpicture}
\begin{axis}[eink, xlabel={$\beta$}, ylabel={$\sigma_{\text{env}}(\beta)$},
  xmin=0, xmax=15, ymin=0, ymax=0.42]
\addplot[black, line width=1.4pt, domain=0.05:15, samples=200]
  {sqrt(x+1)/((x+2)*sqrt(x+3))};
\addlegendentry{$\sigma_{\text{env}}=\dfrac{\sqrt{\beta+1}}{(\beta+2)\sqrt{\beta+3}}$}
\addplot[only marks, mark=*, mark size=2pt, black] coordinates {(2.236,0.1856)};
\addlegendentry{media medida: $0{,}186$}
\end{axis}
\end{tikzpicture}
\figcap{La anchura efectiva de la envolvente es la desviación típica de una
$\mathrm{Beta}(1,\beta+1)$, que es el perfil $(1-r)^\beta$ leído como
distribución de probabilidad. Con el $\beta$ medio de los runs sale $0{,}186$,
es decir $W=2s\sigma_{\text{env}}\approx 0{,}37\,s$ frente a un radio de
soporte $s$: \textbf{el radio sobreestima el tamaño útil casi tres veces}. Ese
factor es exactamente el que separa un $f_1$ autocalibrado que funciona de uno
que nace apagado.}
\end{figure}

\subsection{Qué forma eligió el optimizador}

\begin{figure}[H]\centering
\begin{tikzpicture}
\begin{axis}[eink, height=5.2cm, xlabel={iteración}, ylabel={coeficiente medio},
  xmin=0, xmax=31000, ymin=0, ymax=0.15,
  scaled x ticks=false, scaled y ticks=false,
  ytick={0,0.05,0.1,0.15}, yticklabels={0,{0,05},{0,10},{0,15}},
  xtick={2500,7000,15000,20000,25000,30000},
  xticklabels={{2,5k},7k,15k,20k,25k,30k}]
\addplot[black, line width=1.5pt, mark=*, mark size=1.6pt] coordinates {
  (2500,0.1047) (7000,0.1304) (15000,0.1254) (20000,0.1204) (25000,0.1173) (30000,0.1153)};
\addlegendentry{$a_1$}
\addplot[black, dashed, mark=square*, mark size=1.6pt] coordinates {
  (2500,0.0442) (7000,0.0785) (15000,0.1009) (20000,0.1046) (25000,0.1050) (30000,0.1039)};
\addlegendentry{$a_2$}
\addplot[black, dotted, line width=1.1pt, mark=triangle*, mark size=2pt] coordinates {
  (2500,0.0371) (7000,0.0672) (15000,0.1017) (20000,0.1139) (25000,0.1204) (30000,0.1237)};
\addlegendentry{$a_3$}
\end{axis}
\end{tikzpicture}
\figcap{Los coeficientes medios del run\,7 a lo largo del entrenamiento.
$a_1$ sube, se da la vuelta hacia la iteración 7.000 y baja; $a_3$ sube sin
parar y \textbf{cruza a $a_1$} entre las iteraciones 20.000 y 25.000. El
resultado es un espectro \emph{creciente}: lo contrario de lo que produce una
envolvente, y la firma que aparece en los cuatro runs del átomo sin excepción
(en el run\,5, con la cuarta parte de frecuencia, $a_1$ pasó de $0{,}1598$ a
$0{,}1243$ y $a_3$ de $0{,}0549$ a $0{,}0874$: el mismo movimiento, sin llegar
a cruzarse). Es la señal sin explicar que queda, y la que sugiere la siguiente
palanca: forzar $a_1>a_2>a_3$.}
\end{figure}

% ===========================================================================
\section{El reparto Python / CUDA}

Esta es la decisión de ingeniería con más consecuencias prácticas del
rediseño, y no se ve en las ecuaciones.

\subsection{El tensor de cinco columnas}

El rasterizer recibe, por Gaussiana, un bloque de \textbf{5 floats}
(\code{gabor\_kernel.h:45}, \code{GABOR\_STRIDE 5}):
\[
\bigl[\,a_1,\;a_2,\;a_3,\;\varphi,\;b\,\bigr]
\]
$\varphi$ y $b$ \textbf{no son parámetros entrenables}. Son funciones
diferenciables de $a$ y de las escalas, calculadas en Python
(\code{get\_gabor}, \code{gaussian\_model.py:352--374}). CUDA las trata como
entradas independientes, devuelve $\partial L/\partial\varphi$ y
$\partial L/\partial b$, y es \textbf{autograd} quien compone la cadena hacia
\code{\_a} y \code{\_scaling}.

Las tres razones, en orden de importancia:

\begin{enumerate}
\item \textbf{$\gamma$ y $f_1$ no viven en CUDA.} Barrerlos no recompila nada.
      Todo el sweep de \code{GABOR\_F1} (runs 6 y 7) se hizo con una variable
      de entorno.
\item \textbf{El gradiente ``crece y ganarás textura'' llega solo.} Como
      $\varphi$ depende de $s_u$ por \eqref{eq:phi}, autograd propaga
      $\partial L/\partial\varphi \to \partial L/\partial s_u$
      automáticamente. Ese camino \emph{no existe} en el rasterizer: es puro
      Python.
\item \textbf{Una sola fórmula.} El pedestal se escribe una vez.
\end{enumerate}

El coste: cambiar $a$ de $(N,3)$ a $(N,5)$ \textbf{obliga a recompilar el
rasterizer}. Es el único cambio del rediseño que lo exige.

\subsection{Dónde vive cada símbolo}

\begin{center}\small
\begin{tabular}{@{}lll@{}}
\toprule
\textbf{símbolo} & \textbf{qué es} & \textbf{dónde} \\
\midrule
$a_1,a_2,a_3$ & entrenable & Python + CUDA \\
$\beta$       & entrenable & Python + CUDA \\
$s_u,s_v$     & entrenable & Python + CUDA \\
$\varphi$     & derivado   & Python \\
$b$           & derivado   & Python \\
$\gamma$      & constante  & Python \\
$f_1$         & constante  & Python \\
$\kappa$      & constante  & Python \\
\bottomrule
\end{tabular}
\end{center}

Solo tres familias de parámetros del kernel se entrenan: $a$, $\beta$ y, por
la vía indirecta de $\varphi$, las escalas.

% ===========================================================================
\section{Fase 1 --- Preprocess}

El kernel \code{preprocessCUDA} corre \textbf{una vez por Gaussiana} (no por
píxel) y hace: culling de frustum, matriz de transformación, radio en pantalla,
tiles tocados, color desde armónicos esféricos.

Del bloque Gabor solo hace una cosa (\code{forward.cu:193--196}):

\begin{lstlisting}[language=C]
#pragma unroll
for (int _k = 0; _k < GABOR_STRIDE; _k++)
    a_out[GABOR_STRIDE*idx + _k] = a_in[GABOR_STRIDE*idx + _k];
\end{lstlisting}

Copiarlo al buffer de geometría, igual que $\beta$. Y esto tiene una
consecuencia que conviene tener presente:

\begin{quote}
\textbf{El binning no sabe nada del kernel.} El radio en pantalla, la caja
envolvente y el conjunto de tiles se calculan \emph{solo} con la geometría
($\mu$, escalas, rotación). La onda no ensancha ni estrecha la huella.
\end{quote}

Es correcto, porque el soporte lo fija la envolvente $(1-r)^\beta$, que se
anula en $r=1$ pase lo que pase con $S$. Pero significa que el coste de
memoria del rediseño está aquí: el buffer de geometría pasa de $3$ a $5$ floats
por splat. Con $4{,}9$ M splats son $\sim 39$ MB extra, despreciable frente a
lo que ocupa el modelo.

% ===========================================================================
\section{Fase 2 --- Forward}

Aquí es donde se evalúa el kernel: \textbf{una vez por pareja (píxel, splat)},
que es el bucle más caliente de todo el sistema.

\subsection{Antes del kernel: la geometría}

Cada hilo intersecta el rayo del píxel con el plano del surfel y obtiene las
coordenadas locales $s=(u,v)$ (\code{forward.cu:418--434}):
\begin{align}
\rho_{3d} &= u^2+v^2 && \text{geometría real} \\
\rho_{2d} &= \text{FilterInvSquare}\,\|\,\mathbf{d}\,\|^2 && \text{filtro paso-bajo} \\
\rho      &= \min(\rho_{3d},\rho_{2d}), \qquad r=\sqrt{\rho}
\end{align}
El \code{min} es del 2DGS original: cuando el splat se hace sub-píxel, quien
manda es el filtro de pantalla, no la geometría. Y hay un guard,
\code{if (r2 >= 1.0f) continue}, que garantiza $r<1$ en todo lo que sigue ---
un dato del que dependen varias fórmulas del backward.

\subsection{El interruptor de escala}

\begin{lstlisting}[language=C]
const bool modulate = (rho3d <= rho2d);
const float t = (gabor_mode == GABOR_MODE_DIR) ? s.x : r;
const float S = modulate
              ? gabor_wave(a1,a2,a3,b,phi,t,dS_dt,dS_dphi)
              : 1.0f;
\end{lstlisting}

\textbf{La onda solo se aplica donde manda la geometría real.} Si el splat es
sub-píxel, $S$ se fuerza a $1$ y el kernel es tent puro. El razonamiento:

\begin{itemize}
\item Un splat sub-píxel \emph{no puede mostrar textura}. Modularlo sería
      pintar detalle por debajo de la frecuencia de Nyquist de la pantalla:
      aliasing, no señal.
\item Mantiene intacto el gradiente de la rama $\rho_{2d}$, que es puramente
      geométrico. Sin el interruptor aparecerían términos cruzados entre la
      onda y el filtro paso-bajo.
\end{itemize}

Este interruptor es también el origen del diagnóstico \code{[GABOR-W]}: es la
frontera entre ``la onda está encendida'' y ``este splat es un tent''.

\subsection{La evaluación}

\begin{lstlisting}[language=C]
float dE_dr;
const float E = gabor_envelope(r, beta_j, dE_dr);
kernel = fmaxf(E * S, 0.0f);   // S<0 -> el lobulo negativo no pinta
float alpha = min(0.99f, opa * kernel);
if (alpha < 1.0f/255.0f) continue;
\end{lstlisting}

Con la envolvente (\code{gabor\_kernel.h:70--76}):
\[
E=\bigl(\max(1-r,\,10^{-6})\bigr)^{\beta},\qquad
\frac{\partial E}{\partial r}=-\beta\,(1-r)^{\beta-1}
\]
El piso $10^{-6}$ evita $\mathrm{pow}(0,\beta)$ con $\beta<1$ en el borde.

Y la onda (\code{gabor\_kernel.h:55--65}), que calcula de una vez el valor y
las dos derivadas que hará falta más tarde:
\[
S=b+\sum_n a_n\cos(w_n t),\qquad w_n=(2n-1)\varphi
\]
\[
\frac{\partial S}{\partial t}=-\sum_n a_n w_n\sin(w_n t),\qquad
\frac{\partial S}{\partial \varphi}=-t\sum_n a_n(2n-1)\sin(w_n t)
\]

Tres observaciones sobre el \code{fmaxf(E*S, 0)}:

\begin{enumerate}
\item La cota \eqref{eq:cota} ya garantiza $S\ge0$, así que \emph{en teoría} el
      \code{fmaxf} es código muerto. Está por lo mismo que los demás pisos del
      rasterizer: un \code{.ply} de un run viejo, o un epsilon de float32,
      no deben poder producir opacidad negativa.
\item A diferencia del legacy, aquí \textbf{no hay división}. El legacy dividía
      por $f_0$, y esa división fue la singularidad que reventó el run\,1
      (52,6\,M splats y NaN en la iteración 4.100).
\item \textbf{$K$ no está acotado por 1}, y conviene saberlo. El máximo de $S$
      está en $t=0$ (los tres cosenos en fase) y vale
      \[
      S(0)=b+\ssum=1+\gamma\,\ssum ,
      \]
      así que $K(0)=1+\gamma\ssum>1$: con los coeficientes medios del run\,7
      son $1{,}103$. El pico del kernel \emph{supera} a la tent en un
      10\,\%. No rompe nada --- el \code{min(0.99f, ...)} de $\alpha$ es el
      guard, heredado del 3DGS --- pero significa que
      \textbf{$\ssum$ y la opacidad están acoplados}: subir la amplitud de la
      onda sube también la opacidad del centro, así que parte del gradiente
      que recibe $a$ es ``quiero más opacidad'', no ``quiero más textura''.
      Es una interacción que sale de las ecuaciones y no está medida.
\end{enumerate}

\subsection{Coste}

Por pareja píxel--splat modulada: \textbf{3 senos y 3 cosenos} más un
\code{powf}. El tent solo pagaba el \code{powf}. En unidades de la SFU de la
GPU eso es del orden de $6\times$ más trabajo trascendente en el punto más
caliente del pipeline.

Dos atenuantes medidos: (a) el interruptor de escala deja fuera a los splats
sub-píxel, que en estas escenas son mayoría; (b) los senos y cosenos se evalúan
en \code{sinf}/\code{cosf} (precisión simple, hardware), no en las versiones
de doble precisión.

Hay una optimización pendiente que \emph{no} es del kernel pero le afecta: el
\code{cutoff} sigue en $3{,}0$ (valor gaussiano) cuando el soporte es compacto
en $r=1$, así que se están evaluando del orden de $9\times$ más píxeles de los
necesarios. Ver el punto 3 de la lista de pendientes del proyecto.

% ===========================================================================
\section{Fase 3 --- Backward}

El backward recorre las mismas parejas píxel--splat en orden inverso,
reconstruyendo la transmitancia $T$ hacia atrás.

\subsection{Regla número uno: recomputar, no recordar}

El backward \textbf{vuelve a evaluar el kernel entero} en vez de guardarlo. Es
lo correcto en memoria (guardar $K$ por pareja píxel--splat sería un buffer
prohibitivo) pero implica que las dos fórmulas tienen que coincidir
exactamente. Este proyecto ya pagó una vez ese precio: en los runs 60 y 62 el
clamp de $\alpha$ del backward no coincidía con el del forward y los gradientes
se inflaban en todo el frente del píxel.

De ahí que la definición viva en \textbf{un único header},
\code{gabor\_kernel.h}, que incluyen los dos ficheros. Y de ahí la línea
explícita (\code{backward.cu:368}):

\begin{lstlisting}[language=C]
float alpha = min(0.99f, opa * kernel);   // el MISMO clamp que el forward
const float G = kernel;                   // sin clampar: dalpha/dopa
\end{lstlisting}

\subsection{El punto de partida: $\partial L/\partial\alpha$}

Antes de tocar el kernel, el backward acumula en \code{dL\_dalpha} todo lo que
depende de $\alpha$: color, profundidad, normales, mapa de acumulación y
distorsión (\code{backward.cu:382--436}). Al final multiplica por $T$.

Desde ahí salen \textbf{cinco} gradientes nuevos, más los que ya existían.

\subsection{Los cinco gradientes del bloque Gabor}

Todos salen de derivar \eqref{eq:kernel} directamente, sin división:

\begin{align}
\frac{\partial\alpha}{\partial\beta}   &= \text{opa}\cdot E\cdot S\cdot\ln(1-r) \label{eq:gbeta}\\
\frac{\partial\alpha}{\partial a_n}    &= \text{opa}\cdot E\cdot\cos(w_n t) \label{eq:ga}\\
\frac{\partial\alpha}{\partial b}      &= \text{opa}\cdot E \label{eq:gb}\\
\frac{\partial\alpha}{\partial \varphi}&= \text{opa}\cdot E\cdot\frac{\partial S}{\partial\varphi} \label{eq:gphi}
\end{align}

Compárese con el legacy, cuyo gradiente de $a_n$ era
\[
\frac{\partial\alpha}{\partial a_n}
=\frac{\text{opa}\cdot\beta}{f_0}\,g^{\beta-1}\,(c_n-g),
\]
con un $f_0$ en el denominador, un $g^{\beta-1}$ que hay que apuntalar cuando
$\beta<1$, y una resta $(c_n-g)$ que acopla los tres coeficientes entre sí.
El del átomo, \eqref{eq:ga}, es un coseno multiplicado por dos factores
positivos. \textbf{La simplificación del backward es, en sí misma, un
resultado del rediseño}: menos casos singulares que apuntalar.

La implementación (\code{backward.cu:466--490}):

\begin{lstlisting}[language=C]
const float opaE = opa * E_env;
if (beta_j >= 0.1f) {
    float grad_beta = dL_dalpha * opaE * S_mod
                    * logf(fmaxf(1.0f - r, 1e-6f));
    grad_beta = fminf(fmaxf(grad_beta, -1e-3f), 1e-3f);
    atomicAdd(&dL_dbeta[global_id], grad_beta);
}
if (modulate) {
    const float base = dL_dalpha * opaE;
    ga1 = clamp(base * cosf(w1*t_mod), -50, 50);
    ...
    gph = clamp(base * dS_dphi, -50, 50);
    gb  = clamp(base,           -50, 50);
    atomicAdd(&dL_da[global_id*GABOR_STRIDE + k], g_k);
}
\end{lstlisting}

\subsection{Por qué los clamps son distintos}

$\pm10^{-3}$ para $\beta$ y $\pm50$ para el resto. No es arbitrario:

\begin{itemize}
\item El de $\beta$ es un \textbf{limitador de señal} heredado del kernel beta
      original, donde $\ln g$ podía explotar. En el átomo el logaritmo es
      $\ln(1-r)$, \emph{finito en toda la huella} porque el guard $r<1$ lo
      garantiza --- ya no haría falta, y se mantiene por prudencia.
\item El de $a$, $\varphi$ y $b$ es un \textbf{backstop de NaN}, no un
      limitador: esos gradientes son $O(1)$ y $\pm50$ no los toca nunca. Es
      importante no confundirlos: clampar $a$ a $10^{-3}$ aplastaría la señal
      y el kernel no aprendería nada.
\end{itemize}

\subsection{El gate $\beta\ge0{,}1$ ya no cubre todo}

En el legacy, los splats con $\beta<0{,}1$ se saltaban el gradiente de $\beta$
para evitar NaN. En el átomo ese gate sigue aplicándose \emph{solo a $\beta$};
$a$, $\varphi$ y $b$ quedan siempre vivos, porque \eqref{eq:ga}--\eqref{eq:gphi}
no tienen ninguna singularidad. Un splat con $\beta$ minúsculo puede seguir
aprendiendo forma.

\subsection{La bifurcación radial / direccional en la cadena de $\rho$}

Esta es la parte del backward donde es más fácil equivocarse. La derivada
respecto de la geometría (\code{backward.cu:551--564}):

\begin{lstlisting}[language=C]
float dalpha_dr = opa * dE_dr * S_mod;
if (modulate && gabor_mode == GABOR_MODE_RADIAL)
    dalpha_dr += opa * E_env * dS_dt;
d_alpha_d_rho = dalpha_dr / (2.0f * r_safe);
if (modulate && gabor_mode == GABOR_MODE_DIR)
    dalpha_dsx_extra = opa * E_env * dS_dt;
\end{lstlisting}

El razonamiento:

\begin{itemize}
\item En \textbf{radial} la onda depende de $r$, así que su contribución entra
      entera por la regla de la cadena de $\rho$:
      $\partial\alpha/\partial r = \text{opa}\,(E'S+ES')$, y luego
      $\partial r/\partial\rho = 1/(2r)$.
\item En \textbf{dir} la onda depende de $u=s_x$, que \emph{no pasa por
      $\rho$}. Su término se lleva aparte y se inyecta directamente en la
      componente $x$ del gradiente respecto de las coordenadas locales
      (\code{backward.cu:593--596}):
      \[
      \text{dL\_ds}.x = \underbrace{\text{dL\_d\_rho}\cdot 2 s_x}_{\text{geometría}}
      + \underbrace{\text{dL\_dalpha}\cdot\text{dalpha\_dsx\_extra}}_{\text{onda}}
      + \text{dL\_dz}\cdot T_{w,x}
      \]
\end{itemize}

Si se olvidara ese término, la onda direccional sería invisible para la
geometría: el modelo no podría mover ni rotar el surfel para alinear las
franjas con la textura de la escena, que es justamente lo que se le pide.
Nótese la asimetría: la componente $y$ no lleva término de onda, porque en modo
\code{dir} la onda no depende de $v$.

\subsection{El guard $r_{\text{safe}}$}

$\partial r/\partial\rho = 1/(2r)$ diverge en $r\to0$. El código usa
$r_{\text{safe}}=\max(r,10^{-6})$. La divergencia no es real: en el centro
$\partial\alpha/\partial r\to0$ (tanto $E'$ como $S'$ contienen factores que se
anulan), así que el cociente tiene límite finito; el guard solo evita el
$0/0$ numérico.

\subsection{Una asimetría real que conviene conocer}

El término de fondo se suma a \code{dL\_dalpha} \emph{después} de que se hayan
escrito los gradientes de $a$, $\varphi$, $b$ y $\beta$
(\code{backward.cu:568--571}):
\[
\text{dL\_dalpha} \mathrel{+}= \Bigl(\frac{-T_{\text{final}}}{1-\alpha}\Bigr)
\sum_c \text{bg}_c\cdot \frac{\partial L}{\partial \text{pix}_c}
\]
mientras que \code{dL\_d\_rho} sí lo incluye (es el ``FIX \#3'' documentado en
el código). Es decir: los gradientes del kernel omiten la contribución del
color de fondo.

\textbf{En los runs de este proyecto el término es exactamente cero}, porque
\code{white\_background} está en \code{False} y el fondo es negro
(\code{train.py:79}), así que $\text{bg}\cdot\partial L/\partial\text{pix}=0$.
Pero deja de serlo si alguna vez se entrena con \code{--white\_background}
(por ejemplo, en los datasets sintéticos tipo Blender): ahí $\partial
L/\partial a_n$ quedaría sesgado respecto de la derivada exacta, mientras que
$\partial L/\partial\rho$ no. Es una trampa latente, no un bug activo.

% ===========================================================================
\section{Fase 4 --- Python: restricciones y paso del optimizador}

Cada iteración, después de \code{optimizer.step()}, se aplican tres
correcciones en este orden (\code{train.py:403--456}):

\begin{enumerate}
\item \textbf{Clamp de $\beta$} a $[-4,\,2]$ en el parámetro crudo. El techo
      vive en una sola constante, \code{BETA\_RAW\_MIN/MAX}, desde que un print
      desincronizado anunció durante varios runs un experimento que no se
      estaba corriendo.
\item \textbf{Clamp de la escala} sobre el parámetro crudo, no sobre la
      activación. Mata el ``trinquete'': con \code{torch.clamp} el gradiente no
      pasa por encima del máximo, así que un splat que cruce el techo se
      congelaría para siempre.
\item \textbf{Proyección de $a$} sobre el conjunto factible.
\end{enumerate}

\subsection{La proyección: gradiente proyectado exacto}

\code{project\_a\_} (\code{gaussian\_model.py:404--437}) implementa la
proyección euclídea \emph{exacta} sobre
\[
\mathcal{A}=\{a\in\mathbb{R}^3:\ a_n\ge0,\ \ssum\le S_{\max}\},
\]
que es convexo (una caja intersecada con un semiespacio):

\begin{enumerate}
\item $w=\max(a,0)$ proyecta sobre la caja. Si $\sum w\le S_{\max}$, listo:
      por ser óptimo sobre un conjunto \emph{mayor} que $\mathcal{A}$, también
      es la proyección sobre $\mathcal{A}$.
\item Si no, la restricción de suma está activa y la proyección es la del
      símplex $\{a\ge0,\ \ssum=S_{\max}\}$: ordenar descendente, encontrar el
      mayor $\rho$ con $u_\rho-(\text{cumsum}_\rho-S_{\max})/\rho>0$, restar el
      umbral $\theta$ a todos (Duchi et al., 2008).
\end{enumerate}

Esto no es un clamp heurístico: es \textbf{descenso por gradiente
proyectado}, con las garantías que eso conlleva. En tres dimensiones el
\code{sort} es trivial.

Detalle no cosmético: la tolerancia \code{\_A\_SUM\_TOL} $=10^{-6}$. Sin ella,
como la propia proyección deja $\ssum$ exactamente en la frontera, en float32
una fracción de las filas queda un epsilon por encima y se lanzaría el
\code{sort} sobre millones de filas cada iteración sin cambiar nada.

\subsection{La doble imposición de $a_n\ge0$}

$a_n\ge0$ se impone en \emph{dos} sitios: en \code{project\_a\_} (sobre el
parámetro) y en \code{get\_a} con un \code{clamp(min=0)} (sobre el forward).
El segundo existe porque \code{render.py}, \code{metrics.py} y el visor
\textbf{no pasan por el bucle de entrenamiento}. Los dos tienen que moverse
juntos o aparece inconsistencia train$\leftrightarrow$render. La cota de la
suma, en cambio, no se repite en el forward: como la proyección escribe sobre
el parámetro, lo que se guarda en el \code{.ply} ya es factible.

\subsection{La autocalibración de $f_1$}

Con \code{GABOR\_F1} sin definir, $f_1$ se elige una sola vez en
\code{create\_from\_pcd} para que el interruptor $f\!\cdot\!W=1$ caiga en el
percentil 90 de las escalas iniciales: el 10\,\% de surfels más grandes
desarrolla textura y el resto degenera al tent
(\code{gaussian\_model.py:376--402}).

La anchura efectiva no es el radio del soporte, sino
\begin{equation}
W = 2\,s\,\sigma_{\text{env}}(\beta),\qquad
\sigma_{\text{env}}(\beta)=\frac{\sqrt{\beta+1}}{(\beta+2)\sqrt{\beta+3}}
\label{eq:sigma}
\end{equation}
que es la desviación típica de la distribución $\mathrm{Beta}(1,\beta+1)$, cuya
densidad es proporcional a $(1-x)^\beta$ --- es decir, exactamente el perfil de
la envolvente leído como distribución de probabilidad. Con el $\beta\approx
2{,}24$ medido en los runs sale $\sigma_{\text{env}}\approx0{,}186$, o sea
$W\approx0{,}37\,s$: \textbf{el radio sobreestima el tamaño útil unas tres
veces}. Usar el radio en vez de \eqref{eq:sigma} daría un $f_1$ demasiado bajo
y la onda nacería apagada.

De ahí:
\[
f_1 = \frac{\text{extent}}{2\,s_{p90}\,\sigma_{\text{env}}}
\]
En \emph{flowers} salieron $46{,}6247$ rad/extent.

\begin{quote}
\textbf{Trampa operativa.} \code{render.py} llama a \code{load\_ply}, no a
\code{create\_from\_pcd}, así que \textbf{la autocalibración no corre en el
render} y $f_1$ se queda en $0$. Hay que pasarle el valor por variable de
entorno. El primer render del run\,5 se hizo sin ella: el modelo se evaluó con
\code{legacy/norm/f1=0} y con una cota $\ssum$ más estrecha que la del
entrenamiento, así que \code{load\_ply} \emph{reproyectó el 29,2\,\% de los
splats y mutó el modelo}. Resultado: $-5{,}07$ dB y dianas en el césped, con el
modelo intacto. Hoy \code{render\_server.sh} lee la configuración del log del
train y la verifica.
\end{quote}

\subsection{Tasa de aprendizaje}

$a$ se optimiza con \code{lr = 0.001} (\code{gaussian\_model.py:511}), tres
órdenes por debajo de la de posición. No se ha barrido.

% ===========================================================================
\section{Fase 5 --- Densificación y ciclo de vida}

\subsection{Herencia en clone y split}

$a$ se hereda \textbf{tal cual} en las dos rutas de creación
(\code{gaussian\_model.py:900} y \code{:926}):

\begin{lstlisting}[language=Python]
new_a = self._a[selected_pts_mask].repeat(N, 1)   # split
new_a = self._a[selected_pts_mask]                # clone
\end{lstlisting}

El criterio es el mismo que arregló el trinquete de $\beta$: la escala y la
opacidad son magnitudes \textbf{extensivas} y se reparten entre los hijos
($/0{,}8N$ y $/N$); $a$ y $\beta$ son parámetros de \textbf{forma},
adimensionales, y no se reparten. Repartir $\beta$ entre hijos era exactamente
el bug que degradaba el kernel a caja tras seis generaciones de split.

\subsection{Pero la frecuencia sí se hereda mal (y esto es un hallazgo)}

Aquí hay una interacción que no está escrita en ningún sitio del código y que
explica la principal patología medida. En el split:
\[
s^{\text{hijo}} = \frac{s^{\text{padre}}}{0{,}8N},
\qquad\text{y por \eqref{eq:phi}}\qquad
\varphi^{\text{hijo}} = \frac{\varphi^{\text{padre}}}{0{,}8N}
\]
Como $a$ se hereda intacto pero $\varphi$ cae con la escala, \textbf{cada
split apaga un poco más la onda}. El hijo conserva la \emph{amplitud} de la
modulación del padre pero pierde su \emph{frecuencia efectiva}: ya no le cabe
ni medio ciclo dentro de la huella, y $S\to b\approx$ constante, o sea tent.

Esto es exactamente la deriva medida en los tres runs:

\begin{center}\small
\begin{tabular}{@{}lrrr@{}}
\toprule
apagados ($f\!\cdot\!W<0{,}5$) & run5 & run6 & run7 \\
\midrule
@2500  & 88,6\,\% & 73,1\,\% & 43,3\,\% \\
@7000  & 94,3\,\% & 86,0\,\% & 61,0\,\% \\
@30000 & 93,4\,\% & 83,8\,\% & 64,1\,\% \\
\bottomrule
\end{tabular}
\end{center}

El salto empeora siempre entre 2500 y 7000, que es la ventana de densificación
más intensa. Y explica por qué subir $f_1$ desplaza la curva entera pero no
elimina la pendiente: $f_1$ es una constante fijada con las escalas
\emph{iniciales}, y la densificación mueve el blanco. La única variante viva de
esa línea es \textbf{recalibrar $f_1$ durante el run}.

\subsection{Serialización}

$a$ se guarda en el \code{.ply} como \code{a\_0}, \code{a\_1}, \code{a\_2}
(solo los tres entrenables: $\varphi$ y $b$ se recalculan). Al cargar
(\code{gaussian\_model.py:653--680}) hay dos salvaguardas: si el \code{.ply} es
antiguo y no tiene campos \code{a\_*}, se cae al init del modo; y siempre se
llama a \code{project\_a\_}, avisando \emph{solo} si la violación es de
magnitud real, para no dar falsos positivos por el redondeo del roundtrip.

Ese aviso, \code{[A] load\_ply: ... PROYECTADOS}, es el que delata el problema
de configuración descrito antes.

% ===========================================================================
\section{Verificación}

\code{scripts/test\_gabor\_atomo.py} comprueba tres cosas contra el
rasterizador \emph{real}, no contra una reimplementación:

\begin{enumerate}
\item \textbf{El punto neutro.} Con $a=0$, los modos \code{radial} y \code{dir}
      dan la misma imagen y coinciden con $(1-r)^\beta$ bit a bit.
\item \textbf{Los gradientes} de $a_n$, $\varphi$, $b$, $\beta$ y las escalas,
      contra diferencias finitas: medianas de error del 0,2 al
      2,9\,\%.
\item \textbf{No regresión}: el modo legacy sigue dando lo de antes, con
      2,5\,\% en \code{scaling} como control.
\end{enumerate}

Los porcentajes parecen altos para una comprobación de gradientes, pero son los
esperables en float32 con un rasterizador con early-outs discontinuos
(\code{alpha < 1/255}, \code{T < 1e-4}): un perturbación finita puede cruzar un
umbral y cambiar el conjunto de splats que contribuyen.

% ===========================================================================
\section{Qué se midió}

\subsection{El rediseño funciona; la onda no compra}

\begin{center}\small
\begin{tabular}{@{}lrrr@{}}
\toprule
\textbf{run} & \textbf{PSNR} & \textbf{SSIM} & \textbf{LPIPS} \\
\midrule
run2 legacy      & 20,1602 & 0,5537 & 0,3900 \\
run67 tent       & 20,6684 & 0,5811 & 0,3675 \\
\midrule
run5 átomo $1\times$ & 20,6898 & 0,5759 & 0,3724 \\
run6 átomo $2\times$ & 20,5999 & 0,5763 & 0,3716 \\
run7 átomo $4\times$ & 20,7167 & 0,5758 & 0,3720 \\
\bottomrule
\end{tabular}
\end{center}

Dos lecturas, las dos importantes:

\begin{itemize}
\item Contra el \textbf{legacy}, el átomo gana en las tres métricas
      ($+0{,}53$ dB). El rediseño hace lo que prometía: \emph{aprender la forma
      ya no resta}.
\item Contra el \textbf{tent}, empate dentro del ruido. Y el sweep de $f_1$
      demostró que no es porque la onda esté apagada: con $f_1$ a $4\times$ se
      encendieron $5{,}9$ veces más splats, con más amplitud
      ($\ssum$ medio $0{,}2909\to0{,}3429$) y menos kernels planos
      ($21{,}4 \to 12{,}9$\,\%), y las métricas no se movieron.
\end{itemize}

Conclusión: en \emph{flowers}, \textbf{todo el beneficio del átomo viene de la
envolvente} (el pedestal que hace del tent el punto neutro), no de la
oscilación. Lo cual tiene sentido: el césped y el follaje son textura
estocástica de alta frecuencia sin estructura oscilatoria coherente, y un
armónico aprendido no tiene ahí nada que encajar. Está pendiente comprobarlo en
escenas con estructura (bonsai, kitchen).

\subsection{La firma que sí se repite}

En los cuatro runs del átomo, sin excepción: $a_1$ baja y $a_3$ sube durante el
entrenamiento, hasta que en el run\,7 se cruzan ($a_3=0{,}1237 > a_1=0{,}1153$).
El kernel gasta el presupuesto en el armónico más alto --- espectro creciente,
justo lo contrario de una envolvente. Es la señal con más contenido que queda
sin explicar, y la que apunta a la siguiente palanca: forzar $a_1>a_2>a_3$.

% ===========================================================================
\section{Invariantes}

Lo que no se puede romper sin romper el kernel:

\begin{enumerate}
\item \textbf{Una sola definición.} $K$ y sus derivadas viven en
      \code{gabor\_kernel.h}. Forward y backward lo incluyen los dos.
\item \textbf{$a=0$ debe dar la tent exacta.} Es el control experimental. Si
      un cambio lo rompe, el baseline deja de ser comparable.
\item \textbf{El clamp de $\alpha$ es el mismo en los dos pases.}
\item \textbf{Las variables de entorno del render deben coincidir con las del
      train}, $f_1$ incluido. \code{render\_server.sh} lo verifica solo.
\item \textbf{Cambiar el layout del bloque $a$ obliga a recompilar} el
      rasterizer. Cambiar $\gamma$, $f_1$ o $\kappa$ no.
\item \textbf{Los diagnósticos se leen por histograma}, nunca por el máximo.
      El proyecto tiene ya varios casos de un dial que parecía activo y era un
      no-op.
\end{enumerate}

% ===========================================================================
\appendix
\section{Símbolos}

\begin{center}\small
\begin{tabular}{@{}ll@{}}
\toprule
$r$ & distancia normalizada, $r=\sqrt{\rho}$, $r<1$ \\
$u,v$ & coordenadas locales del surfel \\
$\rho_{3d}$ & $u^2+v^2$, geometría real \\
$\rho_{2d}$ & filtro paso-bajo de pantalla \\
$E$ & envolvente $(1-r)^\beta$ \\
$S$ & onda, ecuación \eqref{eq:onda} \\
$b$ & pedestal, ecuación \eqref{eq:pedestal} \\
$\varphi$ & frecuencia base efectiva \\
$w_n$ & $(2n-1)\varphi$ \\
$t$ & argumento de la onda: $r$ o $u$ \\
$\gamma$ & suelo del pedestal (0,3) \\
$f_1$ & frecuencia en rad por unidad de extent \\
$\kappa$ & dial de contraste (1,0) \\
$\beta$ & exponente de la envolvente \\
$W$ & anchura efectiva, ecuación \eqref{eq:sigma} \\
\bottomrule
\end{tabular}
\end{center}

\section{Mapa del código}

\begin{center}\small
\begin{tabular}{@{}ll@{}}
\toprule
kernel y derivadas & \code{gabor\_kernel.h} \\
copia por-Gaussiana & \code{forward.cu:193} \\
evaluación & \code{forward.cu:474--503} \\
recomputación & \code{backward.cu:333--369} \\
5 gradientes & \code{backward.cu:466--490} \\
cadena de $\rho$ & \code{backward.cu:551--564} \\
$\varphi$ y $b$ & \code{gaussian\_model.py:352} \\
autocalibración & \code{gaussian\_model.py:376} \\
proyección & \code{gaussian\_model.py:404} \\
init & \code{gaussian\_model.py:298} \\
herencia & \code{gaussian\_model.py:900,926} \\
paso del optimizador & \code{train.py:403--456} \\
diagnósticos & \code{train.py:542--600} \\
verificación & \code{scripts/test\_gabor\_atomo.py} \\
\bottomrule
\end{tabular}
\end{center}

\section{Documentos relacionados}

\begin{itemize}
\item \code{docs/rediseno\_kernel\_gabor\_adagar.html} --- el rediseño con
      las 14 figuras y los Monte Carlo del conjunto factible.
\item \code{docs/analisis\_run2\_y\_codigo\_completo.html} --- el diagnóstico
      del kernel legacy.
\item \code{docs/render\_y\_env\_vars.html} --- el pipeline de render y la
      tabla de variables de entorno.
\item \code{CLAUDE.md} --- estado actual, baselines y palancas abiertas.
\end{itemize}

\end{document}
